
Soient @@M0@@ un sous-schéma fermé de @@M1@@, et @@M2@@ l'ouvert complémentaire de @@M3@@ dans @@M4@@. On note @@M5@@ la restriction de @@M6@@ au-dessus de @@M7@@, et @@M8@@ la restriction de @@M9@@ au-dessus de @@M10@@. Enfin, on note encore @@M11@@ les images des @@M12@@ dans @@M13@@ et @@M14@@, et on note :
@@M15@@
les homomorphismes définis de la même manière que @@M16@@.

On considère le diagramme
@@M17@@
où @@M18@@ et @@M19@@ désignent les immersions de @@M20@@ et @@M21@@ dans @@M22@@, @@M23@@ et @@M24@@ les immersions de @@M25@@ et @@M26@@ dans @@M27@@. Montrons que si @@M28@@ et @@M29@@ sont surjectives, il en est alors de même pour @@M30@@.

Les deux lignes horizontales du diagramme sont exactes d'après 1.1. On voit alors que pour montrer que @@M31@@ est surjective, il suffit de montrer que les deux carrés sont commutatifs.

Soit @@M32@@~; on a :
@@M33@@
@@M34@@
la dernière égalité résultant de la formule de projection IV 2.11.1.2. L'égalité @@M35@@ résulte donc de l'égalité @@M36@@, vraie d'après IV 3.1.1, car @@M37@@ est
propre et @@M38@@ plat.

